\section{ML Experiments}

\subsection{Benchmarking OCR model on shipment packing list data}
In logistics domain, shipment packing lists are essential documents that detail the contents of a shipment. Currently, the information from the packing list is manually entered into the system which is time-consuming. To automate this, we experimented with various OCR models to extract text from scanned packing list images.

\subsubsection{Problems in Packing List text extraction}
\paragraph{Inconsistent Data Format}
Packing lists vary significantly in format, font styles, and layouts across different vendors. Some issues we observed in collected data is:
\begin{itemize}
    \item Different layouts, fonts and structures
    \item Tables spanning multiple pages
    \item Handwritten notes and stamps overlapping printed text
    \item Low-quality scans with noise and skew
    \item Varied languages and character sets
    \item Complex tables with merged cells and inconsistent alignments
\end{itemize}

\subsection{Solution \#1: Text extraction through LLM-models}
We experimented with LLM-based models like GPT-4 and Claude to extract structured data from packing list images. The workflow involved:
\begin{enumerate}
    \item Upload the scanned packing list image + prompt.
    \item Use LLM to extract raw text from the image.
    \item Use LLM to parse the extracted text into structured JSON format.
    \item Post-process the JSON to ensure data consistency.
\end{enumerate} 
\paragraph{Assumptions about this approach}
This is not a traditional OCR approach. We are leveraging the LLM's ability to understand context and structure in text. So we have to assume
\begin{itemize}
    \item This is the most viable approach given the business needs of quick turnaround and minimal manual intervention.
    \item The cost of using LLMs is justified by the time saved in manual data entry.
    \item The LLMs can recognize and interpret various formats and layouts at an acceptable accuracy level.
    \item This approach is only to validate the idea before investing in building a custom OCR solution.
\end{itemize}


